\section{Estado del Arte}\label{sec:eda2}

El Informe 1 revisó la literatura sobre complejidad y métodos exactos del MCP (Branch \& Bound, formulaciones ILP y su relajación de Motzkin-Straus). Esta sección se enfoca en los enfoques metaheurísticos, que constituyen el ámbito de este segundo informe.

\subsection{Búsqueda local y metaheurísticas para el MCP}

La dificultad práctica de los métodos exactos en instancias grandes o densas —evidenciada en el Informe 1 mediante la brecha de integralidad del 87.1\,\% en \textit{keller4} y la incapacidad de GLPK de cerrar el árbol de búsqueda en 300\,s— motivó desde temprano el desarrollo de heurísticas de búsqueda local para el MCP. Las primeras propuestas son heurísticas constructivas voraces (\textit{greedy}), del tipo \textit{best-in}/\textit{worst-out}, que construyen o podan un clique vértice a vértice \parencite{pardalos1994maximum}. Estas heurísticas son rápidas pero quedan atrapadas en óptimos locales tan pronto como el clique construido se vuelve maximal, lo que ha motivado dos líneas de mejora principales:

\begin{enumerate}
    \item \textbf{Búsqueda Tabú (Tabu Search).} Introducida por \textcite{glover1989tabu1, glover1990tabu2} y sistematizada en \textcite{glover1997tabu}, la Búsqueda Tabú evita ciclar entre óptimos locales prohibiendo temporalmente revertir movimientos recientes mediante una \emph{lista tabú}, permitiendo excepciones cuando un criterio de \emph{aspiración} se satisface (por ejemplo, si el movimiento conduce a una solución mejor que la mejor conocida). \textcite{battiti2001reactive} aplicaron esta idea directamente al MCP mediante \emph{Reactive Local Search}: alternan entre agregar vértices (mientras el clique no sea maximal) y realizar movimientos de \textsc{Drop}/\textsc{Swap} controlados por una lista tabú cuya duración (\emph{tenure}) se adapta dinámicamente según la frecuencia de repetición de soluciones —una extensión conocida como \emph{Reactive Tabu Search} \parencite{battiti1994reactive}—, reportando resultados competitivos en el benchmark DIMACS con tiempos de cómputo muy inferiores a los métodos exactos.
    \item \textbf{Búsqueda local dinámica y basada en penalización.} \textcite{pullan2006dynamic} proponen \emph{Dynamic Local Search} (DLS-MC), que penaliza dinámicamente los vértices no seleccionados para diversificar la búsqueda sin necesidad de una lista tabú explícita, alternando fases de expansión (\textsc{Add}) y de \emph{plateau search} (movimientos que no cambian $f(S)$ pero sí el conjunto de candidatos, análogos al operador \textsc{Swap} de la Tarea 1). En una línea similar, \textcite{katayama2005effective} combinan una búsqueda de vecindad variable con una estrategia de diversificación basada en el historial de vértices visitados, mejorando la robustez frente a instancias con muchos óptimos locales de calidad similar. \textcite{singh2014_simple} proponen heurísticas más simples basadas en el grado de los vértices dentro de conjuntos independientes mínimos, priorizando velocidad de cómputo sobre robustez.
\end{enumerate}

Un elemento común a estos trabajos —y que se retoma directamente en este informe— es que todos operan sobre una representación de solución que es, por construcción, siempre un clique válido (a diferencia de representaciones binarias penalizadas, más habituales en algoritmos genéticos), evitando así el costo de reparar soluciones infactibles en cada iteración. Esto coincide con el diseño de representación y vecindad Add/Drop/Swap especificado en la Tarea 1 (Sección~\ref{sec:met2}).

\subsection{Diversificación mediante reinicio: Búsqueda Local Iterada}

Un riesgo compartido por las búsquedas locales de trayectoria única (incluida la Búsqueda Tabú) es el estancamiento prolongado en una región del espacio de búsqueda cuando la lista tabú por sí sola no basta para escapar de un conjunto de óptimos locales mutuamente cercanos. El marco de \emph{Búsqueda Local Iterada} (\textit{Iterated Local Search}, ILS) de \textcite{lourenco2003ils} aborda este problema intercalando, tras cada fase de búsqueda local, una \emph{perturbación} de la solución incumbente (t\'ipicamente la eliminación aleatoria de una fracción de la solución) seguida de una reconstrucción, lo que permite explorar una nueva región del espacio sin perder por completo el progreso acumulado. Este informe combina ambas ideas —lista tabú para la intensificación local y perturbación tipo ILS para la diversificación global— en un único algoritmo, descrito en la Sección~\ref{sec:met2}.

\subsection{Contextos de aplicación}

Los dos contextos de aplicación del MCP identificados en el Informe 1 —diseño de códigos correctores de errores mediante la conjetura de Keller \parencite{lagarias1992keller}, directamente vinculada a la familia de instancias \textit{keller} utilizada en este trabajo, y la búsqueda de subestructuras moleculares comunes en bioinformática \parencite{gardiner1997clique}— cobran especial relevancia en el contexto metaheurístico: ambos dominios generan instancias reales cuyo tamaño ($n$ en el orden de cientos a miles de vértices, grafos de producto molecular densos) hace inviable la resolución exacta en tiempo práctico, siendo precisamente el nicho donde las metaheurísticas de búsqueda local como la evaluada en este informe resultan indispensables.
